% 来源：OEBPS/Text/Section0001_split_001.xhtml#sigil_toc_id_10

\BookChapter{2}
\section[第2章\q 简单7步投资法实现财务自由：创造终身收入]{第2章\\ 简单7步投资法实现财务自由：创造终身收入}

\begin{BookQuote}
千里之行，始于足下。

\hfill ——老子
\end{BookQuote}

告诉我，你有没有过这种经历，你知道的……和孩子玩视频游戏，你输得很惨、被羞辱的经历。谁总赢呢？当然是孩子啦！但是孩子是怎么做到的呢？是因为孩子头脑更聪明、动作更快、实力更强吗？

我们来分析下事情的经过。你去侄女或者侄子家做客，孩子见了你会说：“快来和我一起玩游戏吧，好叔叔！”

你马上反对：“不，不，我不会玩这个游戏。你玩你的。”

他们会说：“来吧，很容易的！我一教你就会了。”然后他们射死几个屏幕上跳出来的坏蛋。你仍然不愿意玩，于是，他们就开始恳求你了。“玩吧！玩吧！求你了，求你了，求你了！”你很喜欢这个小家伙，所以只好让步了。然后孩子简单的一句话就让你知道你上套了：“你先。”

于是你下定决心，一定要赢！你要在孩子面前露两手。结果如何呢？砰！砰！砰！只要三四秒，屏幕上的坏蛋就把你打死了。你头部中枪。弹孔冒着烟。

轮到孩子玩了。孩子拿起枪，突然砰、砰、砰一阵扫射。很多坏蛋从天而降，从各个角落嗖嗖地飞来飞去。孩子预料到坏蛋的一举一动，逐个击落——45分钟后，终于轮到你玩第二轮了。

现在你发怒了，更加坚定了非赢不可的信念。这一次，你整整坚持了5秒钟才死掉，孩子上来又一口气玩了45分钟。你终于明白这是怎么回事了。你就是被孩子拉过来当炮灰的。

那么，为什么孩子总能赢？他们的脑子反应更快吗？他们的两只手动作更快吗？不是！是因为他们以前玩过这个游戏很多次。

他们已经掌握了人生获得财富和成功的最大秘诀：他们能预测到前方的路。

记住这句话：预测是终极力量！失败者事后才做出反应，领导者事前就能做出预测。接下来，精英中的精英会教你学习预测的技能，瑞·达利欧、保罗·都铎·琼斯，以及其他50位能够预测前路的金融界领袖人物，他们组成的明星军团知道前方的路。他们来指导你在通往财务自由之路上遇到困难和挑战时，让你一路前行免受伤害。正如瑞·达利欧所说的，前方是一片丛林，布满了能在财务自由之路上杀死你的危险的东西，你需要信得过的向导帮你穿越丛林。有了他们的帮助，我们可以为你制订规划，帮你预测前行的困难和挑战，让你避免不必要的压力，最终到达你理想的财务目的地。

我要简要地概述一下我们要去哪里和本书的基本结构，这样你就可以有效地利用本书了。但在此之前，我们先要搞清楚我们真正的目的。本书致力于达成的最重要的结果是：让你掌握自己的财务，能够拥有一份终身投资收益，再也不用为赚钱去工作，实现真正的财务自由！好消息是，每个人都能实现真正的财务自由。即使在起步时你深陷债务泥潭，毫不夸张地说，只要多花一点儿时间，持续保持专注，运用正确的投资策略，你就能在几年之内获得财务上的安全，甚至是财务自由。

在我们开始按照简单7步投资法一步一步地开始投资之旅之前，让我们先探讨一下，为什么在财务上达到安全无忧过去看起来非常简单，现在却相当困难？其中发生了什么改变？我们需要做什么来应对？我们先来上一堂历史课吧。

\begin{BookQuote}
年轻时没有钱可以，但年老时没有钱绝对不行。

\hfill ——田纳西·威廉斯
\end{BookQuote}

现在，你的财务生活好像比过去困难很多，不是吗？我敢肯定，你一定很困惑，现在想要存点钱养老，然后舒舒服服地享受退休生活，为什么会如此困难。我们已经把退休养老看作社会给予的一项福利：人生神圣不可侵犯的阶段。但是别忘了，退休其实是个相当新的概念，事实上，只有最近一两代人才有退休这个概念——对于我们大多数人来说，只有我们的父母和祖父母两代人，才享受到了退休的福利待遇。在他们之前，人们一般要工作到动不了，不能工作为止。

所有人都是一直工作到死。

还记得历史吗？美国的社会保障制度是什么时候才开始有的呢？那是大萧条期间，时任美国总统富兰克林·德拉诺·罗斯福所创立的，那时并没有照顾老人和病人的社会保障体系。在此之前，“老人”的概念和现在完全不同。那时，美国人的平均预期寿命只有62岁。寿命只有这么长！社会保障体系提供的退休福利，政策规定人要活到65岁才能开始享受，所以并不是每个人都能享受到退休福利，就是享受到了也享受不了多久。事实上，创立社会保障体系的罗斯福总统本人，寿命也没有长到能够享受自己的退休金（尽管他可能并不需要这笔钱）。他死时才63岁。

《美国社会保障法》在大萧条时期减轻了数百万名美国人的痛苦，但是其创立的目的并非取代退休储蓄——只是作为满足最基本生活需求的一个补充而已。而且当时设计的制度也不是为了适应我们今天生活的世界。

世界变化很大，新的现实情况是：已婚夫妇二人其中至少有一位会活到92岁的可能性是50\%；其中一位活到97岁的可能性是25\%。

天哪！我们的平均寿命正在快速接近100岁。

随着寿命延长，我们领取养老金的预期年限会更长——长出很多。50年前，人们领取养老金的年限平均为12年，而如今一个人65岁退休后，预期会活到85岁甚至更久。这样，领取养老金的年限会长达20多年。这还只是平均年限，许多人会活得更长，甚至能领取30年的养老金！

\begin{BookQuote}
用30年工作所交纳的养老金，去负担30年的退休生活，这根本不现实。你不能指望只用前30年每月交纳的10\%的工资收入作为养老金，来支撑后30年的退休生活支出。

\hfill ——约翰·沙文，斯坦福大学经济学教授
\end{BookQuote}

你预计自己能活到多少岁呢？我们看到，各种各样医疗技术的突破创新都会延长我们的寿命好几年，甚至几十年。从干细胞技术到3D打印器官，再到细胞再生，各种技术爆发性涌现。你会在本书第七部分，了解更多最新科技的进步。寿命延长肯定是好事，但是你做好准备了吗？我们很多人都没有做好。

美国万通金融集团相关人员做了一项调查，他们访问了1946—1964年婴儿潮一代出生的人：“你害怕的事情是什么？”

你认为答案会是什么？死亡？恐怖袭击？瘟疫？

错！婴儿潮一代出生的人最害怕的事情是：人活着，钱花没了！

（顺便说一下，死亡，只排在第二位。）

婴儿潮一代出生的人有理由害怕自己的储蓄撑不到去世，千禧一代出生的人也是如此。安永会计师事务所的一项研究表明，75\%的美国人愿意在临死之前看到自己的资产完全消失。社会保障安全网络，如果能够存活到下一代的话，本身并不能提供合理的生活水平。目前，美国人每月能够领取的养老金平均是1294美元。如果你生活在纽约、洛杉矶、芝加哥或迈阿密等大城市，你觉得每个月这点儿钱能撑多久呢？或者如果你住在伦敦、罗马、东京、中国香港、悉尼或新德里等国际大都市，同样是社会保险体系，同样是每月只领取养老金，这又能让你撑上多久呢？无论你生活在哪里，如果没有另一份收入来源，你最终只能老了也得打工，也许做个站在沃尔玛超市门口迎接顾客的最佳着装接待员，一直到你死。

显然，我们需要增加退休后的收入，帮助自己维持更长的退休生活，比以前任何时代都要长——我们现在正处在经济不景气的时期，许多人都在苦苦地挣扎，想要恢复失去的收入水平。

活得更久了，花得更多了，但是经济不景气下收入下降了，我们的经济状况越来越窘迫。我们是如何应对的呢？很多人发现，退休养老这个问题太令人头疼了，简直让人崩溃，于是我们干脆阻止自己去想它，盼望问题会自行消失。根据雇员福利研究所的调查研究，美国在职雇员中，48\%的人从来没有计算过自己退休后总共需要花多少钱养老。是的，48\%！这是个惊人的数字。接近一半的美国人需要从第一步做起，开始制订自己未来的财务规划——算账的时间到了。

那么，解决办法是什么？从第一步开始：做出你一生最重要的财务决策。到你读完本书的时候，你不仅有一个自动的储蓄和投资计划，而且你也知道了不用依靠工作照样能创造收入的方法。

等等！竟有这样的好事！简直太好了，这不可能是真的！你肯定会这样想。任何事情听起来好得简直不能太好了，往往都不是真的，难道不是吗？

但我相信，凡事都有例外。如果我告诉你，现在确实有些金融工具，能让你在市场上涨时赚到钱，在市场下跌时也不会亏一分钱，你会怎么说呢？20年前，对于普通投资者来说，这样的事情根本无法想象。但投资者在2008年市场暴跌一半的行情中，使用这些投资工具，没有亏损一分钱，甚至不会因为担心害怕而失眠。我用这些投资工具让家人获得了财务安全和财务自由。知道自己的收入永远不会被花光，这种感觉简直太棒了。我想确保你和你的家人都有这样的财务安全和财务自由的美好感觉。在本书里，我要展示给你如何创造有保障的终身收入来源。

终身收入，再也不用为了赚钱而被迫工作。

每个月底打开邮箱，原来你看到的只是银行对账单，存款账户里的钱越花越少，你只能幻想存款余额没有减少。相反，创造终身收入之后，月底你的邮箱里出现的不是对账单，而是一张支票，这样的感觉是不是好多了呢？这种好事有办法实现。

在本书第二部分，我会向你展示如何把你的投资打造成赚钱机器，让你稳稳地赚到相当多的养老储备金——我称之为达到“临界数量”。这笔投资让你睡觉的时候也能照样赚钱。用几个简单的策略，你就可以创造出有保障的收入流，让自己能够按照自己的想法来建立、管理、享受自己的私人“养老金”。

也许你很难想象，如今竟然会有这样的投资安排，可以提供给你如下好处：

\begin{itemize}
\item 100\%的本金安全保护，这意味着你绝对不会亏掉本金。
\item 你投资账户的收益与股票市场涨幅（比如标准普尔500指数）直接挂钩。所以，如果股票市场上涨，你就能分享到收益。但是如果市场下跌，你并不会发生亏损。
\item 你也可以把投资账户余额转化成有保障的收入，这些收入你一辈子也花不完。
\end{itemize}

你可以不用想象了，这样的好事就在你眼前。这只是其中的一种，还有许多类似的投资机会，如今像你这样的投资者都可以得到。（你会在第五部分详细地了解到。）

我要明确的是，即使有了终身收入，我也不建议你一到传统的退休年龄就停止工作。当然，可能你也不愿意这样做。研究表明，赚的钱越多，不退休继续工作的可能性就越大。过去，我们的人生目标是发财致富，然后40岁就提前退休。但现在，我们的人生目标是发财致富，然后一直工作到90岁。美国年收入达到7.5万美元（超过50万元人民币）或者更高的人，接近一半的人表示他们永远不会退休，即使退休，至少要工作到70岁。

滚石乐队的主唱米克·贾格尔71岁了，仍然在做世界巡演，像这样工作到老怎么样呢？

或者想想看，企业大亨史蒂夫·韦恩72岁了还在工作。

沃伦·巴菲特84岁了。

鲁伯特·默多克83岁了。

萨默·雷石东91岁了。

尽管他们年纪很大了，钱也非常多了，早就可以退休了，这些企业家仍然经营管理着自己的企业，而且做得很棒。（他们可能以后还会继续工作。）也许你也可以这样，虽然年纪很大了，但仍继续热情工作。

但如果我们不能工作了，或者再也不愿意工作了，会出现什么情况呢？单凭养老金是不能完全支撑我们退休的生活开支的。美国婴儿潮一代出生的人，现在每天都有10000人年满65岁，老年人与年轻人的比例越来越不平衡，至少在我们知道的期限里，是不可能恢复平衡的。1950年，平均16.5个劳动者向社会保障体系交纳的养老金来支撑1个退休者所领取的养老金。2014年，平均2.9个人交纳的养老金来支撑1个人所领取的养老金。

你觉得这样的比例能长期维持下去吗？

《纽约时报》专栏作家、畅销书作家托马斯·弗里德曼，写了一篇文章《现在是401(k)退休计划的世界》。他说：“如果你是自我激励、自力更生的人，哇，太好了，现在这个新世界就是为你量体裁衣的。所有边界都消失了。但如果你不是自我激励、自力更生的人，现在这个新世界会对你构成很大的挑战，因为过去保护你的墙壁、天花板和地板通通消失不见了……限制少多了，但保障也少多了。你做出的贡献大小，会在更大程度上决定你得到福利的多少。作秀不行，你得真正有贡献才行。”

你要是想与父母和祖父母一样，指望着退休后靠着领取养老金过上甜蜜幸福的晚年生活，你还是别做梦了，想想铁匠和电话接线员的命运，以后再也不会有这样的好事了。美国私营企业的员工，总体上只有一半人有退休养老保险，而且其中大部分人现在采用的都是个人完全自理、个人承担所有风险的模式。

如果你是一个市、州或联邦政府的雇员，相当于中国的公务员，那么你还能享受到政府财务所支持的养老金，但随着时间一天天过去，越来越多的人，像底特律、圣贝纳迪诺等那些财政破产的城市里的人，政府雇员心里一直在打鼓，熬到自己退休的时候，这些养老金可能早就没有了。

那么，你的退休养老计划是什么？你有养老金吗？是401(k)退休储蓄计划，还是个人退休金账户呢？如今，有6000万名美国人参加了401(k)退休养老金计划，交纳的养老金规模超过了3.5万亿美元。但如果你加入的是主导市场的高收费投资计划，对于你来说，这样的退休养老金计划的结果会很差劲，甚至亏得很惨。正因如此，如果你参加了401(k)退休养老投资计划，你就必须好好地读第二部分第10章了。你学到的很有用的新东西，让你只要做出简单的改变，就可以完全改变你的人生——带给你平和的心态和你现在需要得到的收益确定性。也就是说：不改变，你老了也不能退休；改变了，你就能提前退休。二者差别很大。

\subsection{死亡与纳税：唯一始终不变的事}

你要小心的敌人是市场波动很快（简直比光速还快）的市场、高得离谱的（和隐性的）费用、陈旧过时的养老保障体系，不要忘记还有我们很好的老朋友——税务员。啊，税务员。你只要赚到钱，税务员就要收走50\%（甚至更多）的所得税。谢谢你为国家做出的贡献。如果你认为隐性费用是你积累财富的唯一拖累因素，那么你错了，你遗漏了最大最严重的拖累因素。

我们都知道纳税会拖累你的财富积累，但在相当大的程度上，很少有人意识到税收拿走我们收入的比例有多大，对我们实现财务自由的能力影响有多大。经验丰富的投资老手非常明白这个道理：重要的不是你赚了多少钱，而是你留到口袋里多少钱。

世界上最伟大的投资者都明白税收效率的重要性。那么随着时间的积累，纳税所造成的破坏性究竟有多大呢？

我们来举个简单的例子：你挣了1美元，你有能力使每年收益翻番，一直持续20年。我们都知道一种投资游戏——复利，对吧？

过了1年，1美元变成了2美元；

过了2年，变成4美元；

过了3年，变成8美元；

过了4年，变成16美元；

过了5年，变成32美元。

你不能用计算器，猜猜看，这样过了20年，你最初的1美元增值到了多少钱呢？

不要作弊。花点时间猜猜看。

通过复利的魔力，20年后，最初的1美元变成了1048576美元！这就是复利令人难以置信的神奇力量。

作为投资者，我们自然想要充分利用复利如此巨大的力量。不过，投资游戏并不是这么简单。在现实世界里，政府首先要你掏钱纳税。税务员正盯着他想分到手的那一块饼呢。在和上面的例子同样的情形下，如果必须交纳所得税，其影响会有多么大呢？还是不要用计算器，先猜猜看。如果你够幸运，每年支付的所得税税率只有33\%，那么你觉得20年之后，你赚到的钱在交完税后，还会剩下多少呢？

不要往后翻，偷看答案，花点时间，认真猜猜看。

好啦！如果税前投资收益是1048576美元……嗯，交纳33\%的所得税后，剩下的投资净收益差不多只有75万美元吗？或者可能更少，只有50万美元吗？再想想。

现在我们接着往下看。在现实世界里，在新的一年开始，复利赚钱机器再次把你的钱翻番之前，我们首先要拿出钱来足额交纳上一年的投资收益所得税，那么税收的吸钱威力，是不是也会像复利的赚钱魔力一样，大到不可思议的程度呢？假设每年投资收益的所得税税率为33\%，同样是经过20年之后，最后你实际获得的税后投资收益只有2.8万多美元。

没错，只有大约2.8万美元！不用每年交纳所得税让你得到的投资收益接近105万美元，而每年交纳33\%的所得税后，你最终拿到手的投资净收益只有2.8万多美元。相差100多万美元——还没算上州税呢！在某些州，比如加利福尼亚州、纽约州和新泽西州，州里征收的所得税也不少。所以你最后拿到手的税后净收益肯定会更低。

当然，在这个例子中，每年收益翻番只是假设，你在现实世界中永远看不到这么高的投资收益率。但是这个简单的例子说明，如果我们忽视考虑税收对财务规划的巨大影响，其后果可能会很严重。

鉴于事情发生在华盛顿州，你认为未来几年的所得税税率会更高还是更低呢？

（你甚至根本不需要回答上述问题。）

在本书第五部分，我会告诉你“投资门路”，直到现在为止能够得到这些投资门路的人，只有经验丰富、手法老到的投资者，或者“超级富豪”。我会让你看看那些最聪明的投资者做了什么——如何避开所得税，《纽约时报》称之为“专给富人的内线秘密”。它就是得到美国国税局批准的避税方法，让你可以免交所得税，从而大大加快财富积累的速度。你不需要是富人，也不需要是名人，普通人都能使用。这种合理避税的方法，基本上能让你的财务增加速度加快25\%\textasciitilde{}50\%，帮助你更快地实现财务独立。不过，具体能快多少取决于你的税级。

\begin{BookQuote}
没有成为自己的主人，就永远不是自由的人。

\hfill ——爱比克泰德
\end{BookQuote}

但是，不管你有没有制订退休养老规划，未来都在飞速到来。根据美国退休养老研究中心的研究，53\%的美国家庭未来“处于危险之中”，因为他们退休后没有足够的收入来继续维持现在的生活水平。一半以上的家庭啊！记住，1/3以上的劳动者所储蓄的养老金还不到1000美元（不包括社会保险养老金和自有住房价值），60\%的劳动者所储蓄的养老金不到25000美元。

怎么会这么少呢？我们不能把这一切都归咎于经济不景气。储蓄危机早在最近这次经济危机之前就开始了。2005年，美国的个人储蓄率为1.5\%；2013年，它是2.2\%（在危机最严重的时候超过了5.5\%）。这是怎么回事呢？我们并不是孤立生活的。我们知道需要储蓄更多，以积累资本做投资。我们为什么没有真正这样做呢？是什么阻碍我们不能更多地储蓄呢？

首先，我们必须承认，人类行为并不总是理性的。有些人花钱买彩票，尽管知道赢得强力球头奖的概率是1/175000000，而我们遭到雷劈，即闪电击中的概率都比它大251倍。事实上，看看下面的统计数据会让你大吃一惊：美国家庭每年花在彩票购买上的钱平均为1000美元。我第一次听到这个统计数据是我的朋友什洛莫·贝纳茨——加利福尼亚州大学洛杉矶分校的著名行为金融学教授——告诉我的。我的第一反应是：“不可能有这么多！”其实，在我最近一次举办的研讨会上，我问现场观众：“买过彩票的人请举手。”结果现场5000名观众中，举手的还不到50人。如果5000人中只有50个人买过彩票，而平均每年购买彩票要支出1000美元，那么这就意味着很多人每年买彩票的实际支出大大高于1000美元。顺便说一下，最高纪录保持者是新加坡，平均每个家庭每年花费4000美元来购买彩票。你能不能大致估计一下，如果把每年买彩票的1000美元、2000美元、3000美元、4000美元拿出来做投资，每年的收益按照复利不断增长，最后价值会增长到多少呢？在下一章中，你会发现只需要很少的本金，通过投资理财使本金不断增值，就能把自己的养老金积累到50万\textasciitilde{}100万美元，甚至更多，而且几乎不需要投入时间管理投资。

现在，我们把讨论的话题转到了行为经济学上，让我们寻找一些小技巧，它们能够造成贫穷与富有的巨大差别。行为经济学家努力想要弄清楚，为什么我们会在财务上犯下错误，如何无意识地轻易改掉这些错误。听起来很酷，对吧？

丹·艾瑞里是杜克大学著名的行为经济学教授。他研究我们的大脑如何经常愚弄我们自己。人类的进化依赖于视觉，人脑的大部分都致力于处理视觉信息。但我们的眼睛经常误导我们。你如果不相信，就看看下面的例子。下图是两张桌子。

\begin{figure}[htbp]
\centering
\includegraphics[width=0.72\linewidth,keepaspectratio]{image/2-1.jpg}
\caption{}\label{fig:2-1}
\end{figure}

如果我问你，哪张桌子更长，是左边窄的桌子，还是右边宽的桌子呢？大部分人会自然而然地选择左边窄的桌子。如果你也给出了同样的答案，那么你就错了。两张桌子的长度一模一样。（你不相信，自己动手量一量吧。）不服，好，我们再来一次。

\begin{figure}[htbp]
\centering
\includegraphics[width=0.72\linewidth,keepaspectratio]{image/2-2.jpg}
\caption{}\label{fig:2-2}
\end{figure}

同样一张图，同样的两张桌子，哪张桌子更长呢？你是不是还觉得左边窄的桌子更长呢？你甚至敢赌上一把。你已经知道答案了，但是你的大脑仍在继续欺骗你。你就是觉得左边窄的桌子看起来更长。你的眼睛所看到的跟不上你的大脑所想到的。“我们的直觉一直在愚弄我们，利用不断重复、事前可预测、始终不变的方式。”艾瑞里在那次令人难忘的TED演讲中说，“对此我们束手无策。”

俗话说眼见为实，事实表明并非如此，我们眼睛看到的东西也会出错。要知道，从理论上讲看东西还是我们人类非常擅长的领域，那么换成我们不擅长的领域，我们犯错的概率会有多大呢？比如财务决策。无论我们认为自己做的财务决策是出色的还是糟糕的，我们内心理所当然地认为自己能够掌控自己所做的决策。然而，科学表明事实并非如此，我们无法掌控自己的决策。

就像上述的视觉错觉一样，我们做决策很容易受到错觉影响。艾瑞里在接受采访时告诉我，我们很多决策错误都要归咎于“认知错觉”。举个比较恰当的例子，你明天要去当地的车辆管理局回答一个关键问题：“你愿意捐赠器官吗？”你认为自己会如何回答呢？有些人会马上回答“愿意”，认为自己很无私、很高尚。另外一些人可能会沉默不语，或者犹豫不决，或者可能会回避这个可怕的问题，拒绝回答。也许你可能会推辞说自己需要时间考虑。但无论你的决策是什么，你认为自己的决策都是自主决策，完全基于个人的自由意志，我的事我自己说了算。你是拥有完全行为能力的成年人，完全可以自主决定是否捐赠个人器官去挽救他人。

但实际情况是另外一回事，你的决策在很大程度上取决于你居住的环境。如果你住在德国，你同意捐赠器官的概率是1/8，即约12\%的人会同意捐赠器官。而德国的邻国奥地利，有99\%的人愿意捐赠自己的器官。瑞典人愿意捐赠器官的比例是8\%，丹麦人只有4\%。这到底是怎么回事，为什么会有这么大的差异呢？

是宗教因素，还是恐惧因素？是文化差异吗？结果表明这些都不是原因。器官捐赠概率的巨大差异，与你个人或文化因素没有一点儿关系。这完全取决于车辆管理局表格上的话是怎么写的。

在器官捐赠概率最低的国家，如丹麦，表格上有个选择框，上面写的是“如果你愿意参加器官捐赠，在此打钩”。而在器官捐赠概率最高的国家，如瑞典，表格上写的是“如果你不想参与器官捐赠，在此打钩”。

这就是秘密所在。没有人喜欢在确认的选择框里打钩。这不是因为我们不愿意捐赠器官，只是人性懒惰，不想费事而已。就是那么一点点的惰性，造成了世界上的差别。

如果碰上的问题非常严重，那么我们容易被吓呆，一动不动，什么也不做，或者想也不想，只是被动地按照别人说的去做。这不是我们的错，我们天生就是这样处理事情的。人们在器官捐赠的问题上很犹豫，并不是因为不关心，而是因为太关心了。做出决定很困难，也很复杂，很多人根本不知道应该做什么。“因为我们不知道应该做什么，所以我们就去依赖别人，别人怎么定我们就怎么做。”艾瑞里说。

同样的惰性，或者说被动地接受别人给我们所选的路，有助于解释为什么只有1/3的美国工人充分利用了现有的退休养老金计划。同样的原因也可以解释为什么只有少数人为自己的未来制订了财务规划。长期的财务规划看起来太复杂了。我们不知道该做什么，所以我们干脆撒手不管，什么都不做。

艾瑞里告诉我，提到有形的物理世界，我们很清楚自身的能力是有限的，于是想办法用设计和建设来克服这些限制。比如，我们建造了台阶、坡道和电梯，让自己走得更远，上得更高。“但出于某种原因，在设计医疗保健体系、退休养老体系和股票市场体系时，我们不知道为什么就忘了我们的能力是有限的。”艾瑞里说，“我认为，如果我们能够认识了解我们大脑认知的局限性，就像我们认识了解自己身体的局限性一样，我们就能把整个世界设计得更美好。”

还记得瑞·达利欧所说的吗？要进入丛林，他会问自己的第一个问题是：“我不知道的有哪些？”如果知道自己的局限，你就能适应调整，进而取得成功。如果不知道自己的局限，你就会受到伤害。

我写本书的目的是要唤醒众人，传授给他们知识和工具，让他们马上可以掌控自己的财务生活。我创建的投资理财规划非常简单易行。为什么？因为投资理财规划，如果用起来太复杂、太难、太花时间，人们就会停止使用。正如我们在车辆管理局的表格上看到的那样，复杂是执行的敌人。正是因为这个原因，我将这个投资理财规划分为简单的7步，而且开发了一个功能强大的手机App，它完全免费，用来一步一步地指导你完成这7个步骤。你可以一边前行一边检查自己的进度，每完成一步就要庆祝一次。这个手机App会支持你，回答你的问题，甚至在需要的时候提醒你一下。因为你会变得兴奋激动，努力追求最美好的目标，可是也会偶尔分心或出现惰性，这会让你偏离目标。设计自动化系统的目的就是要防止你出现这样的情况。猜猜看，你做好之后的结果会如何？你一次做好了，就一劳永逸。你的投资理财规划已经按部就班地执行到位，从此之后，你只要每次花上一小时，一年只需一两次检查一下，确保自己的航向正确就行了。有什么理由不坚定原来的路线继续前行呢？这条路让你能够得到终身的财务安全、财务独立、财务自由，而且它耗时很少，能让你把足够多的时间用在对自己至关重要的事情上。

我希望，读到这里，你的大脑正在激烈地做思想斗争。我知道，到目前为止，你有太多东西需要好好思考，但是我承诺，要给你的财务生活创造出持续性的突破进展，我想让你清楚地了解前面要走的路。所以，让我们一起快速浏览一下通向财务自由的简单7步投资法吧。

如果你属于在博客、推特、微博中长大的一代，我猜你会不耐烦地说：“就简单的7步，这么一点东西，你竟然写了一大本书！你为什么不把这7步总结成一句话或者浓缩在一张图表里呢？”我可以做到。但知道信息是一回事，真正理解信息并遵照执行又是另一回事。不能被贯彻执行的信息，一点儿用也没有。记住，我们被淹没在信息的海洋里，但我们真正渴求的是智慧。

所以我先要改变你的思维，然后才能为后面的每一步做好准备。只有通过这种方式，你才能为采取必要的行动做好准备，才能保证自己按照这七步一步一步地前进，最终实现财务自由。

我撰写本书的目的是让你掌握一个主题，而这个主题正折磨着大多数人，因为他们从来没有花时间掌握其中的基本技能来让自己彻底摆脱痛苦。掌握意味着深入理解——谁都能读完后记住内容，这样便觉得自己学会了其中的东西。但真正要掌握它，你需要经过三个层次。

第一个层次是思维认知上的理解，即你理解概念的能力。这个层次我们每个人都能达到。许多人对个人理财和投资都有一些认知理解。但是你的认知理解，再加上3美元，差不多够给你到星巴克买一杯咖啡。我的意思是信息本身并没有价值，所以认知层面上的理解仅仅是第一步。

想要开始得到真正的价值，你必须进展到第二步，即心理情绪上的控制。一遍又一遍地重复，次数多了，你自然就知道了，这些东西激发出你内心很多的情感——欲望、渴望、恐惧和担忧，现在你能意识到这些东西了，也能持续运用你学到的东西了。

但最终的掌握是身体行动上的掌握，也就是说，你不用想自己做什么，你的行动成了你的第二天性。要想达到身体行动上的掌握，唯一的方法就是不断地重复、重复、再重复。我的那位杰出的人生导师吉姆·罗恩告诉我，重复是技能之母。

我讲一个典型的例子，正是我觉得自己能力比较欠缺的领域——跆拳道。我20岁出头的时候，很想拿到跆拳道黑带（最高等级），我有幸认识了李俊九大师，还成了他的好朋友。就是他把跆拳道带到了美国，他训练指导过李小龙和拳王阿里。我跟李大师说，我想用历史上最短的时间获得黑带，为了打破这一纪录，我什么事情都愿意做，刻苦训练，完全投入，严守纪律。李大师同意和我一起旅行，在旅途中指导我完成训练。训练非常残酷！我首先要工作，举办研讨会，培训讲课，之后再去场馆内参加训练，我凌晨1点赶到，然后李俊九大师指导我训练三四个小时。这样下来，我一天最多只能睡上4个小时。

有天晚上，完全相同的一个动作，我练了很长时间，特别长，至少练了300遍。我实在受不了了，转身问大师：“大师，我们什么时候可以练下一个动作啊？”他看着我，很严厉，回答道：“你这个家伙，跟个蚂蚱一样，一点儿耐心也没有，你现在练的就是下一个动作。其实，你根本分不清你现在练的这个动作和前面练的那个动作间的区别。这就表明你还是个业余爱好者。能不能看出这些细微的差别，就是专业大师与业余爱好者之间的差距所在。要掌握，就需要这种精细的重复。多重复一次，你才能多学一点儿。”

你明白我的意思了吗？我写书时，也有意做了这样的设计安排，让你不能只用一个下午就能轻松地翻完它。

本书完全不同于你以往读的任何一本书，因为它体现了我非常独特的教学风格。我会问你很多问题，有时你看到的一些事实情况和话语，你在前面已经读过了。你还会看到很多感叹号，这可不是编辑错误！我是有意使用这种技巧的，标记出关键思想，把知识牢牢嵌入你的大脑、身体和心理，让你的这种行为变成自动反应。那时，你就会开始看到成果，收获你想要的也是应得的回报。你愿意接受挑战吗？

要记住：这不仅是一本书，它还是一个行动方案。每一章的设计，都是要帮助你从投资理财的角度正确地理解自己的水平，帮助你缩小实际与期望之间的差距。让你来做这些功课，目的是要武装你。这不只是为了你的今天，更是为了未来的人生做好准备。我知道，你会在不同的阶段重新阅读本书，然后把自己再提升到更高的水平。

\subsection{第一部分\q 欢迎来到丛林：旅程从这一步开始}

就像所有伟大的冒险家一样，我们一开始先要为旅程确定方向。在第一部分第4章，你会学习了解到更多关于财富心理学的知识，分析阻碍我们积累财富的心理因素，并学习一些简单的解决之道。你会发现投资真正的目的是什么，从而释放出最佳财务突破策略的巨大力量。然后，在第5章，我们就要像火箭发射一样点火升空了，从此开始你的简单7步投资法。在这里你会迈出第一步，做出一生中最重要的财务决策。这一章是必读内容。你会学到哪怕只用很少一点的钱作为原始本金，只要学会运用复利的神奇力量，你就绝对可以实现终身财务独立，即使你的收入不高，也没有交上好运。要想激活这套投资理财规划，你首先要把省出的一部分收入作为储蓄以积累本金，来投资赚取复利收益。之后你不再是经济上的消费者，而是一个所有者——一个投资者，为未来押注。你会学会如何打造自己的自动赚钱机器，这个系统将会为你创造终身收入，即使在睡觉，你也照样能赚钱。

\subsection{第二部分\q 成为内行：加入之前先了解游戏规则}

也许你听说过这句老话：“有钱人碰到有经验的人，最后的结局是，原来有经验的人带着钱走了，有钱人只带着经验走了。”既然现在你决定要成为一个投资者，本书第二部分会给你解释一些投资游戏的核心规则，好让你不会轻易地输掉比赛，不会成为那些很有经验的游戏高手的猎物。证券市场就是瑞·达利欧所说的投资丛林，路线图会告诉你穿越投资丛林的行进路线，里面所有危险的区域都标上了大大的红X。它们都是那些市场营销人士编造的投资神话——有些人称之为投资谎言。那些投资机构设计安排这些投资谎言是有套路的，目的是一步一步地把你跟你的钱分开。你会了解到为什么基金公司广告上吹嘘的投资业绩并不是你实际能够得到的投资业绩。我知道这听起来不可思议，但是，你认为1\%的管理费就是你所支付的所有成本费用吗？其实，这只是潜在十几项费用中的一种。平均而言，每只公募基金会吃掉你长期潜在投资收益的60\%！记住，只学习第二部分其中的一小部分内容，你就能省下25万\textasciitilde{}45万美元。即使你未来一生的投资收益保持不变，这也能让你的口袋里又多出25万\textasciitilde{}45万美元！这个数据，我并不是随口乱说的，也不是胡乱推算的，这都是基于白纸黑字的学术研究成果。我们会讨论那些推销到期日基金和免佣金基金的投资广告的宣传骗术，传授你一些投资真知来武装自己，保护自己免受那些投资机构的欺骗。它们经常会定制一些产品，专门设计一些投资策略，但是这样做的目的是使投资机构自己的利益最大化，而不是为了使个人投资者的利益最大化。读完本书的第二部分后，你就可以迈出你简单7步投资法的第二步了，即使你只有很少的一点儿本金，你也可以像一个内行一样去做投资了。

\subsection{第三部分\q 梦想的代价，让游戏能赢}

我们来一起探索你的财务梦想，设立一些切实可行的目标，让你确实有可能赢得投资游戏。很多人根本不了解他们究竟需要多少钱，才能获得财务安全、财务独立、财务自由，或者他们脑子里所想象的目标财富金额非常巨大，以至被吓坏了，无从制订一个以此为目标的投资理财规划。但是在第三部分第15章你会弄清楚你真正想要的金钱目标是什么，这个符合实际的目标会让你非常激动，特别是当你意识到你的梦想比你想象的还要近的时候。这不仅是梦想，而且你会把这个梦想变成现实，方法是实践一个投资理财规划，我将在第16章中详述。对于不同的人来说，实现人生梦想的投资理财规划并不相同，我们开发了一个软件，用来为你量身打造专属于你的投资理财规划。你可以通过网页或者手机应用程序定制，做好规划后你可以保存下来，以后想怎么修改都行，直到你找到一个真正切实可行的投资理财规划为止。如果你觉得自己实现梦想的速度不够快，我会在第三部分告诉你利用五种方法来加快速度。等你走完了第三步，你不但知道如何积累财富来支撑自己未来幸福自由的退休生活，也会知道如何一边前进一边享受一路的风景。

\subsection{第四部分\q 做出你一生最重要的投资决策}

现在你已经能像内行一样思考了，你也知道了投资游戏的重要规则，你还学会了如何把这个游戏变成你可以获胜的游戏。现在时间到了，你需要做出一生中最重要的投资决策——该把自己的钱投到哪里呢？每部分投资占所有资金的比例有多大呢？这就是资产配置。每个诺贝尔经济学奖获得者，每个对冲基金经理，每个顶尖的机构投资者，每个专业投资人都告诉我，资产配置是成功投资的关键，但是基本上99\%的美国人只知道一点儿，甚至根本不懂资产配置。为什么？也许是因为资产配置看起来太复杂了。不过在第四部分第22章，我会把资产讲得非常简单易懂，还会向你演示到哪里去找专家在线帮助你做好资产配置。合理的资产配置就好像把你的投资分开放到两大类不同的木桶里。第一类木桶非常安全，能给你带来心灵的平静，而第二类木桶的风险就大得多了，但是可能会有更大的收益增长潜力。这就是终极木桶清单！它是不是非常简单易懂呢？当你完成了第四步后，你不但知道如何从没钱变得很有钱，而且知道如何一直守住财富。

\subsection{第五部分\q 只赚不赔：创造一个终身收入计划}

如果没有钱花，那么投资对你的生活有什么好处呢？大多数人过去一直专注于把更多的钱投入401(k)养老金投资计划，或者投入个人退休金账户里面。时间一长，这就成了习惯，结果他们忘记了未来，他们需要像花掉工资收入一样花掉投资账户里的钱。由于投资账户的余额会变动（记住并非只会上涨！），我们必须制订一个收入计划并且让这个收入计划保持稳定。还记得2008年股市大崩盘吗？如何在下一次遇到股市大崩盘时，保护自己免受重大伤害呢？如何建立一个良好的投资组合，让你可以避免因急欲补偿损失而低价出售同一种证券，结果导致双重损失的结果呢？你怎么知道你不会最终落到大多数人最恐惧的人生结局——钱花没了，人还活着呢？你很长寿，这是上天赐给你的福气，但如果你老了，钱却花光了，你就不会觉得自己有福气了。在第五部分里，我们会给你深刻的洞察力，帮你解开金融界保守得最严密的秘密，帮助你研究设计出一个靠谱的终身收入计划——一份确定的收入源源不断地流进你的账户。此时你的心里就有底了，能够获得财务心理上真正的平静。我们会探索创造性的投资方法，让你能够完全停止亏损或者把亏损限制在非常小的范围之内，同时增加你的收益——使用的投资工具就是银行、大企业和一些世界上最有钱的人所青睐的投资工具。他们知道，而你不知道的投资秘诀是什么呢？那就是如何让你只有赚钱的可能而没有亏钱的可能，而且确保你的收益不被税收吃掉。

\subsection{第六部分\q 像万里挑一的大师那样投资：亿万富翁战术手册}

关于全球经济状况，我们很想知道，未来面临哪些有利因素和挑战——我们如何发展到现在的这种状况？未来全球经济又会发展成什么样子？关于这些问题，我将会介绍世界上思考最透彻而且影响力最大的思想家的分析。然后，你会遇到12位杰出的投资游戏大师，在金融界他们个性最鲜明、成就最卓越，你可以学习12位顶尖投资高手的投资策略，这正是指引你稳稳地取得非凡业绩的投资秘籍。我会问保罗·都铎·琼斯，1987年他是如何预测到美国股市“黑色星期一”大崩盘，从而一个月就获得60\%的收益的；21年后，在整个股市下跌近50\%的大熊市里，在整个世界似乎要再次崩溃的经济环境中，他又是如何逆市获得近30\%的收益的。此外，我们还会分析他是如何避免亏损，如何安然地度过每一种你能够想象得到的市场波动期，一分钱没亏过，连续28年获得盈利的。你在这一部分还会遇到更多的投资大师，比如查尔斯·施瓦布、卡尔·伊坎、托马斯·布恩·皮肯斯、瑞·达利欧、约翰·博格，他们并非出生于富贵人家，只能依靠自己的努力奋斗，逐步从贫穷到富有。他们是如何成为投资的顶尖高手，进而积累了巨大财富的呢？我会问这些顶尖的投资高手，金钱对于他们来说意味着什么。我还会揭开他们实际的投资组合中都有哪些股票。等你走完第6步，你就知道那些万里挑一的投资高手是如何投资的了。

\subsection{第七部分\q 只管去做，尽情享受，尽情分享}

在第七部分，我们会制订一个行动计划，帮助你过上更美好、更充实、更富有、更快乐的生活。我们会讨论如何坚定目标，运用突破性的新技术，让你未来的生活比你想象的更加美好，甚至在不久的将来就能实现，我敢保证这会让你大吃一惊。这一点与大多数人的想法完全相反。美国全国广播公司联合《华尔街日报》进行的一项民意调查表明，76\%的美国人，比例创历史新高，认为孩子未来的生活水平会比自己这一代更差！但你会以内行的眼光来分析我们这个时代最卓越的大脑对未来的分析预测。我们会倾听两位大师的高见，我的朋友雷·库兹韦尔，人称当代爱迪生，还有X奖的创始人彼得·戴曼迪斯，他们会畅谈最新的科技创新：把你的个人电脑变成制造工厂的3D打印机，自动驾驶汽车，帮助截瘫患者重新行走的外骨骼，单细胞长成的假肢，这些科技创新会在很近的将来显著地改变我们的生活。我希望这些能激励你，也让你看到，即使有些地方做得很糟糕，在财务方面的行动没有整合好，你仍然能得到更高的生活品质。对于那些掌握了这些资源，并且能够按照我们制定的投资理财规则采取行动的人而言，你的未来拥有无限的可能、无限的机会、无限的幸福。

我们用一个简单的事实来圆满地结束本章：生活的秘诀就是给予。与他人分享，不仅会大大提高你的生活品质，也会大大提升你的快乐体验。你会学到一些新技术，让你给予的过程既没有痛苦又快乐有趣。当你充实了思想又积累了财富，实现了物质和精神双丰收后，我希望你能做得更好，让自己有充足的能力去帮助他人。别忘了，你现在就是我们的搭档。你在读本书的时候，其实你买书的钱有一部分已经捐出来了，我们一起捐赠餐食给缺少食物的人。

\begin{BookQuote}
我相信，人们寻找的不是生活的意义，而是活着的体验。

\hfill ——约瑟夫·坎贝尔
\end{BookQuote}

以上就是实现财务自由的简单7步投资法，我尽可能讲得通俗易懂。现在该你采取行动了，按照这7步法，一步一步地前进，每次一步，直到完成任务，达成财务自由的终极目标。

你想如何读完整本书呢？用什么方式阅读最适合你呢？我们现在也一起制订一个简单的阅读计划吧。有些人可能会抽出一个周末，坐下来一口气读完整本书。你要是这样做了，那么你就和我一样疯狂一样痴迷，我们真是志同道合！如果你不能抽出整整一个周末的时间，那就换个轻松的方式，比如一天读一章，或一周读一部分。有空就读上几页，几周后你就能把书读完了。不管采用什么方式，只管去读吧。

这是决定你一生的投资理财之旅，这是值得你努力学习掌握的投资理财之旅！来，我们一起出发吧！
